\section{Return to the four-ball}\label{sec:rtc}

We work with the torus cover of \S\ref{sec:fer} and with unit representatives in $\C^3\times\C$.
Write $\mathcal F(Z)=Z_0^3+Z_1^3+Z_2^3$ and fix $0<h<d_n<1/3$. The source set is
\[
 \mathcal N_h=\{(Z,U):|Z|=1,\ \mathcal F(Z)=0,\ \pi|U|^2\le h\}.
\]
The cover map $\mathsf N(u,v,w)=(Z(u,v),w)$ takes $\{p\le h\}$ into $\mathcal N_h$, and the three marked base points are the
zeros of $Z_0$ on the torus; on the cover the zero set is their full $\Z^2$-periodic inverse image
(Lemma~\ref{lem:fer-bpoints}). Two pairs are gauge equivalent if $(Z',U')=(cZ,c^3U)$ with $|c|=1$. We
prove that the transfer $\mathcal T_{\rm u}$ is injective on representatives and gauge equivariant, so that
$\mathcal T_{\rm u}\circ\mathsf N$ identifies exactly the deck-equivalent points.

\subsection{Conventions and the ambient form}\label{sec:rtc-conventions}

For $Z\ne0$ use
\[
 \eta_H=\frac{\operatorname{Im}(Z^*dZ)}{2\pi|Z|^2},\qquad
 \omega_{\rm FS}=d\eta_H=dd^c\log|Z|^2.
\]
For a unit vector put $H_Z=\{\xi:Z^*\xi=0\}$.
On $H_Z$, $\omega_{\rm FS}(\xi_1,\xi_2)=\pi^{-1}\operatorname{Im}(\xi_1^*\xi_2)$: at $|Z|=1$,
$d\eta_H=\frac1{2\pi}\sum_jd\operatorname{Im}(\bar Z_jdZ_j)+(\text{terms containing }d|Z|^2)$, and
$\frac1{2\pi}d\operatorname{Im}(\bar Z_jdZ_j)(\xi_1,\xi_2)=\frac1\pi\operatorname{Im}(\bar\xi_{1,j}\xi_{2,j})$, while
$d|Z|^2=2\operatorname{Re}(Z^*\,\cdot\,)$ vanishes on $H_Z$. All forms below are invariant under the gauge action and annihilate its
generator, hence descend to the quotient; we evaluate them on horizontal representatives.

\begin{definition}[Ambient potential on representatives]\label{def:rtc-ambient}
\lean{SymplecticPacking.Return.PhiE,SymplecticPacking.Return.lamE,SymplecticPacking.Return.OmegaE}
On $(\C^3\setminus0)\times\C$ set
\[
 q(Z,W)=\frac{|W|^2}{|Z|^6},\quad
 \Phi_E^{(\delta)}=\log|Z|^2+\delta\log(1+q),\quad
 \Omega_E^{(\delta)}=dd^c\Phi_E^{(\delta)},\quad \Omega_E=\Omega_E^{(d_n)}.
\]
On unit representatives put $p_W=d_n|W|^2/(1+|W|^2)$ and
$X_T=\{(Z,W):|Z|=1,\ \mathcal F(Z)=TW\}$.
Use $E_3$ as shorthand for these representative data with gauge $(Z,W)\sim(cZ,c^3W)$.
For $W\ne0$, write $\alpha=d\arg W/(2\pi)-3\eta_H$.
\end{definition}


A variable unit phase changes $\eta_H$ by $d\chi/(2\pi)$ and the fibre angular form by $3d\chi/(2\pi)$.
Thus $\alpha$ is gauge invariant and zero on the gauge generator. All calculations can be pulled back along a smooth
family of unit representatives. The potential and its derivatives are explicit functions on the indicated ambient space.

\begin{lemma}[$d^c$ of the norm]\label{lem:rtc-dc}
\lean{SymplecticPacking.Return.dc_log_normW, SymplecticPacking.Return.dc_log_one_add}
\uses{def:rtc-ambient}
On $\{W\ne0\}$, $d^c\log\|W\|^2=\alpha$. On $E_3$,
\[
 d^c\log\big(1+\|W\|^2\big)=\frac{\|W\|^2}{1+\|W\|^2}\,\alpha ,
\]
where the right side extends smoothly across $\{W=0\}$.
\end{lemma}

\begin{proof}
Over an affine chart of the unit representative space choose a holomorphic lift $\hat Z$ of $[Z]$ and the holomorphic fibre coordinate
$\upsilon$ with $[Z,W]=[\hat Z,\upsilon]$. Then $\|W\|^2=|\upsilon|^2\|\hat Z\|^{-6}$, so
$d^c\log\|W\|^2=d\arg\upsilon/2\pi-3\,d^c\log\|\hat Z\|^2$. In the unit representative $Z=\hat Z/\|\hat Z\|$ one has
$W=\upsilon/\|\hat Z\|^3$, so $\arg W=\arg\upsilon$; and since $\operatorname{Re}(\hat Z^\dagger d\hat Z)=\frac12d\|\hat Z\|^2$,
\[
 \eta_H=\frac{Z^\dagger dZ}{2\pi i}=\frac{\operatorname{Im}(\hat Z^\dagger d\hat Z)}{2\pi\|\hat Z\|^2}=d^c\log\|\hat Z\|^2 .
\]
This proves the first identity. The second follows from $d^c\phi(\mu)=\phi'(\mu)\,d^c\mu$ and $d^c\mu=\mu\,d^c\log\mu$ for
$\mu=\|W\|^2$. The right side is smooth because
$\|W\|^2\alpha=\|\hat Z\|^{-6}\big[(\operatorname{Re}\upsilon\,d\operatorname{Im}\upsilon-\operatorname{Im}\upsilon\,d\operatorname{Re}\upsilon)/2\pi
-3|\upsilon|^2d^c\log\|\hat Z\|^2\big]$.
\end{proof}

\begin{lemma}[The horizontal Hermitian blocks]\label{lem:rtc-kahler}
\lean{SymplecticPacking.Return.OmegaE_closed,SymplecticPacking.Return.OmegaE_herm,SymplecticPacking.Return.OmegaE_pos}
At $|Z|=1$, for $(\xi,\varsigma)\in H_Z\times\C$,
\begin{equation}\label{eq:rtc-OmegaE}
 \Omega_E^{(\delta)}((\xi_1,\varsigma_1),(\xi_2,\varsigma_2))
 =\frac1\pi\operatorname{Im}(A_\delta\xi_1^*\xi_2+B_\delta\bar\varsigma_1\varsigma_2),\quad
 A_\delta=1-\frac{3\delta|W|^2}{1+|W|^2},\quad B_\delta=\frac{\delta}{(1+|W|^2)^2}.
\end{equation}
For $\delta=d_n$, $A>1-3d_n>0$ and $B>0$.
The form is closed, and its restriction to every complex horizontal constraint kernel
$\{(\xi,\varsigma):d\mathcal F_Z\xi=T\varsigma\}$ is nondegenerate.
\end{lemma}
\begin{proof}
Write $q=|W|^2|Z|^{-6}$, which is invariant under $(Z,W)\mapsto(\lambda Z,\lambda^3W)$, $\lambda\in\C^*$; hence
$\Omega_E^{(\delta)}$ is the sum of $\omega_{\rm FS}$ and the pullback of a form from the quotient, and it suffices to evaluate on
$H_Z\times\C$ at $|Z|=1$. There $\partial|Z|^2=Z^*dZ$ vanishes, so $\partial q=\bar W\,dW$ and
$\partial\bar\partial q=dW\wedge d\bar W-3|W|^2\sum_jdZ_j\wedge d\bar Z_j$. Hence
\[
 \partial\bar\partial\log(1+q)=\frac{\partial\bar\partial q}{1+q}-\frac{\partial q\wedge\bar\partial q}{(1+q)^2}
 =\frac{dW\wedge d\bar W}{(1+|W|^2)^2}-\frac{3|W|^2}{1+|W|^2}\sum_jdZ_j\wedge d\bar Z_j ,
\]
with no mixed terms. Since $dd^c=\frac i{2\pi}\partial\bar\partial$ and $\frac i{2\pi}dz\wedge d\bar z(\xi_1,\xi_2)=\frac1\pi\operatorname{Im}(\bar\xi_1\xi_2)$, adding
$\omega_{\rm FS}$ gives \eqref{eq:rtc-OmegaE}. For $\delta=d_n$, $A\ge1-3d_n\frac{|W|^2}{1+|W|^2}>1-3d_n>0$ and $B>0$, so the
Hermitian form $A\xi_1^*\xi_2+B\bar\varsigma_1\varsigma_2$ is positive definite; its restriction to any complex subspace is
positive definite, so the imaginary part is nondegenerate there. Closedness is $d\,dd^c=0$.
\end{proof}


%----------------------------------------------------------------------
\subsection{The central normal injection}\label{sec:rtc-injection}

The injection matches the moment $p=\pi|U|^2\in[0,d_n)$ of $N$ with $p_W\in[0,d_n)$ fibrewise.

\begin{construction}[Weight-three injection]\label{cons:rtc-injection}
\lean{SymplecticPacking.Return.injection}
\uses{def:rtc-ambient, def:cub-model}
For a unit pair $(Z,U)$ with $p<d_n$ put
\[
 \iota[Z,U]=[Z,W],\qquad W=\frac{\sqrt\pi\,U}{\sqrt{d_n-p}},\qquad p=\pi|U|^2 .
\]
\end{construction}

\begin{lemma}[The fibre injection]\label{lem:rtc-injection}
\lean{SymplecticPacking.Return.injection_bijective,SymplecticPacking.Return.injection_normW,SymplecticPacking.Return.injection_pW,SymplecticPacking.Return.injection_normW_le}
At each unit $Z$, the radial maps
\[
 U\mapsto\frac{\sqrt\pi\,U}{\sqrt{d_n-\pi|U|^2}},\qquad
 W\mapsto W\sqrt{\frac{d_n}{\pi(1+|W|^2)}}
\]
are smooth inverses from $\{\pi|U|^2<d_n\}$ to $\C$ and back.
They commute with the weight-three gauge action and preserve $Z$.
Moreover $|W|^2=p/(d_n-p)$, $p_W=p$, and on $\mathcal N_h$ one has $|W|^2\le h/(d_n-h)$.
\end{lemma}
\begin{proof}
Substitute each radial expression into the other and use positivity of its denominator.
The norm and moment identities follow by squaring the positive radial factor.
The last bound follows from the increasing function $p\mapsto p/(d_n-p)$.
\end{proof}


\begin{proposition}[The central pullback on the cover]\label{prop:rtc-central}
\lean{SymplecticPacking.Return.injection_pullback}
On $\{p<d_n\}$ in the torus cover,
$(\iota\circ\mathsf N)^*\Omega_E=\Omega_0$.
\end{proposition}
\begin{proof}
By Lemma~\ref{lem:rtc-dc}, on unit representatives
$d^c\Phi_E=\eta_H+d_n\frac{|W|^2}{1+|W|^2}\alpha=\eta_H+p_W\alpha$. Along $\iota$ we have $p_W=p$ and $\arg W=\arg U$
(Lemma~\ref{lem:rtc-injection}), and along $\mathsf N$, $Z^*\eta_H=\eta$ and $U=w$. Hence
$(\iota\circ\mathsf N)^*d^c\Phi_E=\eta+p\bigl(\frac{d\arg w}{2\pi}-3\eta\bigr)=(1-3p)\eta+\frac12\operatorname{Im}(\bar w\,dw)$, which is
smooth across $w=0$, and its exterior derivative is $\Omega_0$ by Lemma~\ref{lem:fer-Omega0}.
\end{proof}


%----------------------------------------------------------------------
\subsection{Transport to $T=1$}\label{sec:rtc-transport}

At a point of $X_T$ the tangent space $TX_T$ is the kernel of the derivative of $\mathcal F(Z)-TW$. The transport field is
the lift of $\partial_T$ which is $\Omega_E$-orthogonal to $TX_T$ and preserves the constraint.

\begin{construction}[Transport field]\label{cons:rtc-transport}
\lean{SymplecticPacking.Return.nvec, SymplecticPacking.Return.Gfun, SymplecticPacking.Return.Nden,
 SymplecticPacking.Return.transportField}
\uses{lem:rtc-kahler}
For $(Z,W,T)\in S^5\times\C\times\R$ put $f=\mathcal F(Z)$, $\bar Z^2=(\bar Z_0^2,\bar Z_1^2,\bar Z_2^2)$,
\[
 n=3\bar Z^2-3\bar fZ,\qquad G=\|n\|^2,\qquad N_{\rm den}=\frac GA+\frac{T^2}B ,
\]
with $A,B$ as in Lemma~\ref{lem:rtc-kahler}. On the open set $\mathcal O=\{N_{\rm den}>0\}$ define
\begin{equation}\label{eq:rtc-transport}
 \dot Z=\frac{W\,n}{A\,N_{\rm den}},\qquad \dot W=-\frac{T\,W}{B\,N_{\rm den}},\qquad \dot T=1 .
\end{equation}
\end{construction}

Since $A,B>0$ everywhere, $N_{\rm den}\ge0$, with equality only where $T=0$ and $n=0$; so $\mathcal O$ is the complement of
$\{T=0,\ n=0\}$, and the field is smooth on $\mathcal O$.

\begin{lemma}[Properties of the transport field]\label{lem:rtc-transport-field}
\lean{SymplecticPacking.Return.n_horizontal, SymplecticPacking.Return.n_dF, SymplecticPacking.Return.G_eq,
 SymplecticPacking.Return.transport_gaugeEquiv, SymplecticPacking.Return.transport_horizontal,
 SymplecticPacking.Return.transport_constraint, SymplecticPacking.Return.transport_normW,
 SymplecticPacking.Return.transport_orthogonal, SymplecticPacking.Return.transport_isotropic}
\uses{cons:rtc-transport, lem:rtc-kahler}
\begin{enumerate}[label=(\roman*)]
\item $Z^\dagger n=0$, $n^\dagger\xi=d\mathcal F_Z(\xi)$ for $\xi\in H_Z$, and $G=9\sum_j|Z_j|^4-9|f|^2$.
\item The field is equivariant under $(Z,W)\mapsto(e^{i\chi}Z,e^{3i\chi}W)$ and satisfies $Z^\dagger\dot Z=0$; hence it is
tangent to $S^5\times\C\times\R$. Its horizontal component is $V=(\dot Z,\dot W)\in H_Z\times\C$.
\item Along the field, $\frac{d}{dT}\mathcal F(Z)=\frac{WG}{AN_{\rm den}}=\frac{d}{dT}(TW)$, so $\mathcal F(Z)-TW$ is constant;
$\|Z\|$ is constant; and $\frac{d}{dT}|W|^2=-\frac{2T|W|^2}{BN_{\rm den}}\le0$ for $T\ge0$.
\item At a point of $X_T\times\{T\}$ in $\mathcal O$, $V$ is $\Omega_E$-orthogonal to $\ker(d\mathcal F_Z,-T)\subset H_Z\times\C$.
\end{enumerate}
\end{lemma}

\begin{proof}
(i) $Z^\dagger(3\bar Z^2)=3\sum_j\bar Z_j^3=3\bar f$, so $n$ is the orthogonal projection of $3\bar Z^2$ to $H_Z$; hence $Z^\dagger n=0$
and $\|n\|^2=\|3\bar Z^2\|^2-|3\bar f|^2$. For $\xi\in H_Z$,
$d\mathcal F_Z(\xi)=\sum_j3Z_j^2\xi_j=(3\bar Z^2)^\dagger\xi=n^\dagger\xi+3fZ^\dagger\xi=n^\dagger\xi$.

(ii) Under the gauge action $f\mapsto e^{3i\chi}f$ and $n\mapsto e^{-2i\chi}n$, while $G,A,B,N_{\rm den}$ are invariant; so
$\dot Z\mapsto e^{i\chi}\dot Z$ and $\dot W\mapsto e^{3i\chi}\dot W$. By (i), $Z^\dagger\dot Z=0$, so $\operatorname{Re}(Z^\dagger\dot Z)=0$ and the
field is tangent to $S^5\times\C\times\R$. The kernel of $T(S^5\times\C)\to TE_3$ is spanned by $(iZ,3iW)$, and $H_Z\times\C$ is a
complement of it, which gives the descent and the representative.

(iii) By (i), $\frac{d}{dT}\mathcal F(Z)=n^\dagger\dot Z=WG/(AN_{\rm den})$, and
$W+T\dot W=W\big(1-\frac{T^2}{BN_{\rm den}}\big)=W\frac{G/A}{N_{\rm den}}$. Further $\frac{d}{dT}\|Z\|^2=2\operatorname{Re}(Z^\dagger\dot Z)=0$ and
$\frac{d}{dT}|W|^2=2\operatorname{Re}(\bar W\dot W)$.

(iv) By (i), in the gauge of Lemma~\ref{lem:rtc-kahler} the tangent space $\ker(d\mathcal F_Z,-T)\subset H_Z\times\C$ is the kernel of the complex linear
functional $\ell(\xi,\varsigma)=n^\dagger\xi-T\varsigma$ on $H_Z\times\C$. For the Hermitian form
$\mathfrak h=A\,\xi_1^\dagger\xi_2+B\,\bar\varsigma_1\varsigma_2$ of \eqref{eq:rtc-OmegaE} and $\nu_\ell=(n/A,-T/B)$ one has
$\mathfrak h(\nu_\ell,\kappa)=\ell(\kappa)$ for all $\kappa$. Since $V=(W/N_{\rm den})\,\nu_\ell$ and $\mathfrak h$ is conjugate linear in the first
slot, $\mathfrak h(V,\kappa)=0$ for every $\kappa\in\ker\ell$, hence $\Omega_E(V,\kappa)=\frac1\pi\operatorname{Im}\mathfrak h(V,\kappa)=0$.
Moreover $\ell(V)=W(G/A+T^2/B)/N_{\rm den}=W$: the field is the unique $\Omega_E$-orthogonal lift of $\partial_T$ along which the
derivative $n^\dagger\dot Z-T\dot W-\dot TW$ of the constraint vanishes.
\end{proof}

\begin{lemma}[Denominator]\label{lem:rtc-denominator}
\lean{SymplecticPacking.Return.Nden_pos, SymplecticPacking.Return.Nden_lower}
\uses{cons:rtc-transport, lem:rtc-transport-field}
Let $\mathcal C=\{(Z,W,T):\|Z\|=1,\ \mathcal F(Z)=TW,\ 0\le T\le1\}$. On $\mathcal C$, $N_{\rm den}>0$: for $T=0$ one has $G\ge3$
and $N_{\rm den}\ge3$, and for $T>0$, $N_{\rm den}\ge T^2/B>0$. If moreover $|W|\le W_{\max}$ and
$T_0=\min\big(1,1/(\sqrt6\max(1,W_{\max}))\big)$, then $N_{\rm den}\ge\min\big(\frac32,\,T_0^2/d_n\big)$.
\end{lemma}

\begin{proof}
By Cauchy--Schwarz, $\sum_j|Z_j|^4\ge\frac13\big(\sum_j|Z_j|^2\big)^2=\frac13$, and $0<A\le1$, so $N_{\rm den}\ge G/A\ge G$. On $\mathcal C$,
$|f|=T|W|$. For $T=0$ this gives $G=9\sum_j|Z_j|^4\ge3$. If $T\le T_0$, then $9|f|^2\le9T_0^2W_{\max}^2\le\frac32$, so $G\ge\frac32$. If
$T\ge T_0$, then $N_{\rm den}\ge T^2/B\ge T_0^2/d_n$, since $B\le d_n$.
\end{proof}

\begin{proposition}[Protected transport]\label{prop:rtc-transport}
\lean{SymplecticPacking.Return.transport_mem,SymplecticPacking.Return.transport_hasDerivAt,SymplecticPacking.Return.transport_smooth,SymplecticPacking.Return.transport_injective,SymplecticPacking.Return.transport_pullback}
Fix $w_*=h/(d_n-h)$ and let $\tau_T$ be the flow, from $T=0$, of the cut field defined in the proof. It is a smooth
injective map, commuting with the gauge action. For every unit pair $(Z,W)\in X_0$ with $|W|^2\le w_*$ the trajectory
solves \eqref{eq:rtc-transport} on $[0,1]$, stays in $X_T\cap\{|W|^2\le w_*\}$, and $\tau_1$ maps
$X_0\cap\{|W|^2\le w_*\}$ into $X_1$ with $\tau_1^*(\Omega_E|_{X_1})=\Omega_E|_{X_0}$ there.
\end{proposition}
\begin{proof}
Let $T_0=T_0(\sqrt{w_*})$ and $c_*=\min(\frac32,T_0^2/d_n)>0$ be the bound of Lemma~\ref{lem:rtc-denominator}. For
$a<b$ write $S_{a,b}(x)=S((x-a)/(b-a))$ and $P_{a,b,c,d}(x)=S_{a,b}(x)(1-S_{c,d}(x))$, and put
\[
 \chi=S_{c_*/2,3c_*/4}(N_{\rm den})\,P_{1/2,3/4,5/4,3/2}(|Z|^2)\,S_{0,1}(w_*+2-|W|^2)\,P_{-1,-1/2,3/2,2}(T),
\]
where off the unit sphere $n$ is replaced by $3\bar Z^2-3\bar f Z/|Z|^2$ and $N_{\rm den}$ is computed with it. The field
$\chi\cdot(\dot Z,\dot W,1)$ (with $\dot T$ replaced by $\chi$) is smooth on $\C^3\times\C\times\R$, because the first factor
vanishes where $N_{\rm den}\le c_*/2$, and it has compact support in $\{\frac12\le|Z|^2\le\frac32,\ |W|^2\le w_*+2,\ |T|\le2\}$. Its
flow is therefore complete, each time-$T$ map is a diffeomorphism (hence injective), and it commutes with the gauge action
because every factor is gauge invariant and the field is equivariant (Lemma~\ref{lem:rtc-transport-field}(ii)).

On $\mathcal C_*=\{|Z|=1,\ \mathcal F(Z)=TW,\ |W|^2\le w_*,\ 0\le T\le1\}$ every factor of $\chi$ equals one: $N_{\rm den}\ge c_*$ by
Lemma~\ref{lem:rtc-denominator} with $W_{\max}=\sqrt{w_*}$. Along the uncut field, $|Z|$ and $\mathcal F(Z)-TW$ are constant and $|W|^2$
is nonincreasing for $T\ge0$ (Lemma~\ref{lem:rtc-transport-field}(iii)). Hence the uncut trajectory from a point of
$X_0\cap\{|W|^2\le w_*\}$ stays in $\mathcal C_*$ while it exists, so it is also a trajectory of the cut field, which is defined for
$T\in[0,1]$; by uniqueness the two coincide on $[0,1]$.

For the pullback, let $x(\sigma)$ be a smooth family of such initial points and $F(\sigma,T)=\tau_T(x(\sigma))$. Since
$\mathcal F(Z)-TW$ vanishes identically along $F$, a tangent vector $(\kappa,t)$ satisfies
$\ell(\kappa)=tW$, where $\ell=(d\mathcal F_Z,-T)$. Replace $\kappa$ by its horizontal representative
$\kappa_H$ modulo the gauge generator $G_{\rm gauge}=(iZ,3iW)$. On $X_T$ one has
$\ell(G_{\rm gauge})=3i(\mathcal F(Z)-TW)=0$, and $\Omega_E$ annihilates this generator, so the constraint and
the pairing with $\Omega_E$ are unchanged. Now $\partial_TF=(V,1)$, and by Lemma~\ref{lem:rtc-transport-field}(iv),
$\ell(V)=W$ and $V$ is $\Omega_E$-orthogonal to the horizontal kernel of $\ell$. Writing
$\kappa_H=tV+(\kappa_H-tV)$ gives $\Omega_E(V,\kappa)=\Omega_E(V,\kappa_H)=0$. Thus
$\iota_{\partial_T}F^*\Omega_E=0$. As $F^*\Omega_E$ is closed, Cartan's formula gives
$\partial_T(F_T^*\Omega_E)=d\,\iota_{\partial_T}F^*\Omega_E=0$, so $F_1^*\Omega_E=F_0^*\Omega_E$, which is the asserted identity.
\end{proof}


%----------------------------------------------------------------------
\subsection{Return to the Fubini--Study form}\label{sec:rtc-return}

On $X_1$ we have $W=\mathcal F(Z)$, so the projection $\mathrm{pr}(Z,W)=Z$ has inverse the graph map
$\mathrm{gr}[Z]=[Z,\mathcal F(Z)]$. Let $u=|\mathcal F(Z)|^2/\|Z\|^6$, a smooth function on unit representatives, and for $s\in\R$
\[
 \omega^{\rm gr}_s=\omega_{\rm FS}+s\,d_n\,dd^c\log(1+u)=\mathrm{gr}^*\Omega_E^{(sd_n)} .
\]
The last identity holds because $\mathrm{gr}^*\|W\|^2=u$; in particular $\omega^{\rm gr}_0=\omega_{\rm FS}$ and
$\omega^{\rm gr}_1=\mathrm{gr}^*\Omega_E=(\mathrm{pr}^{-1})^*(\Omega_E|_{X_1})$.

The return is the Moser flow for the path $\omega^{\rm gr}_s$: since $\partial_s\omega^{\rm gr}_s=d_n\,dd^c\log(1+u)$, we need
$\iota_{Y_s}\omega^{\rm gr}_s=-d_n\,d^c\log(1+u)$. The field $Y_s$ below solves this in closed form because $n$ is an
eigenvector of the Hermitian matrix $M_s$ of $\omega^{\rm gr}_s$; its factor $\frac12$ is the ratio of the factor $\frac1\pi$
of the K\"ahler convention to the factor $\frac1{2\pi}$ of Lemma~\ref{lem:rtc-dc-horizontal}.

\begin{construction}[Return field]\label{cons:rtc-return}
\lean{SymplecticPacking.Return.omegaS, SymplecticPacking.Return.returnField}
\uses{lem:rtc-transport-field}
In unit representatives, with $f$, $n$, $G$ as in Construction~\ref{cons:rtc-transport}, put
\[
 A_s=1-3s\,d_n\frac{u}{1+u},\qquad B_s=\frac{s\,d_n}{(1+u)^2},\qquad M_s=A_sI+B_s\,nn^\dagger\ \text{ on }H_Z,
\]
\begin{equation}\label{eq:rtc-return-field}
 Y_s=-\frac{d_n}{2}\,\frac{f\,n}{(1+u)(A_s+B_sG)}\in H_Z .
\end{equation}
\end{construction}

\begin{lemma}[Graph forms]\label{lem:rtc-graph}
\lean{SymplecticPacking.Return.omegaS_graph, SymplecticPacking.Return.OmegaE_herm,
 SymplecticPacking.Return.As_gt_half, SymplecticPacking.Return.Bs_nonneg, SymplecticPacking.Return.returnDen_gt_half, SymplecticPacking.Return.Ms_nvec}
\uses{cons:rtc-return, lem:rtc-kahler}
For $s\in[0,1]$ and $\xi_1,\xi_2\in H_Z$, $\omega^{\rm gr}_s(\xi_1,\xi_2)=\frac1\pi\operatorname{Im}(\xi_1^\dagger M_s\xi_2)$. Moreover
$u\le1$, $A_s\ge1-\frac32d_n>\frac12$ and $B_s\ge0$; hence $M_s\ge\frac12I$ and every $\omega^{\rm gr}_s$ is positive on $H_Z$.
\end{lemma}

\begin{proof}
By Lemma~\ref{lem:rtc-transport-field}(i), $d\,\mathrm{gr}$ maps the tangent vector represented by $\xi\in H_Z$ to the vector represented by
\[
 (\xi,d\mathcal F_Z(\xi))=(\xi,n^\dagger\xi)
\]
in the gauge of Lemma~\ref{lem:rtc-kahler}. At
$W=\mathcal F(Z)$ one has $|W|^2=u$, so $A_{sd_n}=A_s$ and $B_{sd_n}=B_s$, and \eqref{eq:rtc-OmegaE} gives the Hermitian form
$A_s\xi_1^\dagger\xi_2+B_s\overline{n^\dagger\xi_1}\,n^\dagger\xi_2=\xi_1^\dagger M_s\xi_2$. For $\|Z\|=1$,
$|f|\le\sum_j|Z_j|^3\le\max_j|Z_j|\sum_j|Z_j|^2\le1$, so $u\le1$, $u/(1+u)\le\frac12$ and $A_s\ge1-\frac32d_n>\frac12$ because $d_n<\frac13$.
\end{proof}

Thus, unlike $\Omega_E$ on $E_3$, whose horizontal coefficient is only bounded below by $1-3d_n$, the graph forms are
uniformly non-degenerate as $d_n\to\frac13$, and the denominator $A_s+B_sG$ of \eqref{eq:rtc-return-field} is at least
$\frac12$ for $s\in[0,1]$.

\begin{lemma}[$d^c$ along the graph]\label{lem:rtc-dc-horizontal}
\lean{SymplecticPacking.Return.dc_log_one_add_u}
\uses{cons:rtc-return, lem:rtc-transport-field}
For $\xi\in H_Z$ (unit representative), $d^c\log(1+u)(\xi)=\dfrac{\operatorname{Im}\big(\bar f\,n^\dagger\xi\big)}{2\pi(1+u)}$.
\end{lemma}

\begin{proof}
Where $f\ne0$, as in Lemma~\ref{lem:rtc-dc},
$d^c\log(1+u)=\frac{u}{1+u}\big(d^c\log|\mathcal F|^2-3\,d^c\log\|Z\|^2\big)=\frac{u}{1+u}\big(\frac{d\arg\mathcal F}{2\pi}-3\eta_H\big)$. On
$H_Z$, $\eta_H(\xi)=0$ and $d\arg\mathcal F(\xi)=\operatorname{Im}\big(d\mathcal F_Z(\xi)/f\big)=\operatorname{Im}(n^\dagger\xi/f)$, and $u/f=\bar f$. Both sides
are smooth, so the identity extends to $\{f=0\}$.
\end{proof}

\begin{lemma}[The return field]\label{lem:rtc-return-field}
\lean{SymplecticPacking.Return.returnField_spec, SymplecticPacking.Return.returnField_gauge}
\uses{lem:rtc-graph, lem:rtc-dc-horizontal}
$Y_s$ is smooth in $(s,Z)$ for $s$ in a neighbourhood of $[0,1]$, gauge equivariant and tangent to $S^5$, hence descends to a
time-dependent smooth vector field on unit representatives, and
\[
 \iota_{Y_s}\omega^{\rm gr}_s=-d_n\,d^c\log(1+u)\qquad(s\in[0,1]).
\]
\end{lemma}

\begin{proof}
Under the gauge action $f\mapsto e^{3i\chi}f$ and $n\mapsto e^{-2i\chi}n$, so $Y_s\mapsto e^{i\chi}Y_s$, and $Z^\dagger Y_s=0$. The
denominator is positive for $-\frac1{9d_n}<s<\frac2{3d_n}$, an interval containing $[0,1]$: for $s\ge0$,
$A_s+B_sG\ge A_s\ge1-\frac32s\,d_n$ since $u/(1+u)\le\frac12$; for $s<0$, $A_s\ge1$ and, since
$0\le G\le9\sum_j|Z_j|^4\le9$ and $u\ge0$, $B_sG\ge-9|s|\,d_n$, so $A_s+B_sG\ge1-9|s|\,d_n$. Since
$M_sn=(A_s+B_s\|n\|^2)n=(A_s+B_sG)n$, we have $M_sY_s=-\frac{d_n}2\frac{fn}{1+u}$, and for $\xi\in H_Z$, by Lemma~\ref{lem:rtc-graph}
and Lemma~\ref{lem:rtc-dc-horizontal},
\[
 \omega^{\rm gr}_s(Y_s,\xi)=\frac1\pi\operatorname{Im}\big((M_sY_s)^\dagger\xi\big)
 =-\frac{d_n}{2\pi(1+u)}\operatorname{Im}\big(\bar f\,n^\dagger\xi\big)=-d_n\,d^c\log(1+u)(\xi).\qedhere
\]
\end{proof}

\begin{proposition}[Unit-representative return]\label{prop:rtc-return}
\lean[needed]{SymplecticPacking.Return.returnFlow_pullback,SymplecticPacking.Return.returnFlowInv_pullback,
SymplecticPacking.Return.returnFlow_smooth,SymplecticPacking.Return.returnFlowInv_smooth,
SymplecticPacking.Return.returnFlowInv_returnFlow,SymplecticPacking.Return.returnFlow_returnFlowInv,
SymplecticPacking.Return.nsq_returnFlow,SymplecticPacking.Return.returnFlow_gaugeEquiv}
Let $\psi_s$ be the flow from $s=0$ of the cut field of the proof. It is defined for all $s$, preserves $|Z|=1$, commutes
with the gauge action, and its time-$s$ maps are diffeomorphisms. On the unit sphere it is the flow of $Y_s$ for
$s\in[0,1]$, and $\psi_s^*\omega_s^{\rm gr}=\omega_{\rm FS}$. Consequently $(\psi_1^{-1})^*\omega_{\rm FS}=\omega^{\rm gr}_1$, and
$(\psi_1^{-1}\circ\mathrm{pr})^*\omega_{\rm FS}=\Omega_E|_{X_1}$.
\end{proposition}
\begin{proof}
Put $m=1-\frac32d_n>0$ and let the cut field be
\[
 S_{m/2,3m/4}(D_s)\,P_{1/2,3/4,5/4,3/2}(|Z|^2)\,P_{-1,-1/2,3/2,2}(s)\,Y_s,\qquad D_s=A_s+B_sG,
\]
with $n$ extended off the sphere as in Proposition~\ref{prop:rtc-transport}. It is smooth (the first factor vanishes where
$D_s\le m/2$) and compactly supported, so its flow is complete with diffeomorphic time maps; it is gauge equivariant because $Y_s$
is (Lemma~\ref{lem:rtc-return-field}). On the unit sphere $Y_s$ is horizontal, so $d|Z|^2(Y_s)=2\operatorname{Re}(Z^\dagger Y_s)=0$ and the
sphere is invariant. For $|Z|=1$ and $s\in[0,1]$, $D_s\ge A_s\ge1-\frac32sd_n\ge m$ (Lemma~\ref{lem:rtc-graph}), so every cutoff
factor equals one and the cut field is $Y_s$. Since $\partial_s\omega^{\rm gr}_s=d\bigl(d_n\,d^c\log(1+u)\bigr)$ and
$\iota_{Y_s}\omega^{\rm gr}_s=-d_n\,d^c\log(1+u)$ (Lemma~\ref{lem:rtc-return-field}), Cartan's formula gives
$\frac d{ds}\psi_s^*\omega^{\rm gr}_s=\psi_s^*d\bigl(\iota_{Y_s}\omega^{\rm gr}_s+d_n\,d^c\log(1+u)\bigr)=0$, hence
$\psi_s^*\omega^{\rm gr}_s=\omega^{\rm gr}_0=\omega_{\rm FS}$. Finally $\mathrm{gr}\circ\mathrm{pr}=\mathrm{id}$ on $X_1$, so
$\mathrm{pr}^*\omega^{\rm gr}_1=\mathrm{pr}^*\mathrm{gr}^*\Omega_E=\Omega_E|_{X_1}$.
\end{proof}

%----------------------------------------------------------------------
\subsection{Equivariant injectivity and the ball chart}\label{sec:rtc-avoid}

From now on, $\mathcal T_{\rm u}$ uses the protected cutoff $w_*=h/(d_n-h)$ and is considered on $\mathcal N_h$.
Let $\mathsf r Z=(\zeta Z_0,Z_1,Z_2)$ and $\mathsf r(Z,W)=(\mathsf r Z,W)$.
It preserves the norm, Fermat equation and source moment and commutes with the unit gauge action.

\begin{lemma}[Fixed representatives modulo gauge]\label{lem:rtc-fixed}
\lean{SymplecticPacking.Return.fixed_CP2,SymplecticPacking.Return.fixed_N}
For a nonzero vector, $\mathsf r Z$ is a scalar multiple of $Z$ exactly when $Z_0=0$ or $Z_1=Z_2=0$.
For a central unit pair $(Z,U)$, the equation $\mathsf r(Z,U)=(cZ,c^3U)$ is possible exactly when $Z_0=0$.
\end{lemma}
\begin{proof}
A nonzero last coordinate forces $c=1$ and then $Z_0=0$; otherwise both last coordinates vanish.
The latter alternative is excluded at a central unit vector by $\mathcal F(Z)=0$.
When $Z_0=0$, the pair is fixed with $c=1$.
\end{proof}


\begin{lemma}[Equivariance]\label{lem:rtc-equivariant}
\lean{SymplecticPacking.Return.injection_equivariant, SymplecticPacking.Return.transport_equivariant,
 SymplecticPacking.Return.returnFlow_equivariant, SymplecticPacking.Return.compositeT_equivariant}
\uses{cons:rtc-injection, prop:rtc-transport, prop:rtc-return}
$\mathsf r$ preserves $\mathcal F$, $\omega_{\rm FS}$, $\|W\|$, $\Omega_E$, $u$, each $X_T$ and $\{(Z,U):|Z|=1,\mathcal F(Z)=0,p<d_n\}$, and it maps the fields
\eqref{eq:rtc-transport} and \eqref{eq:rtc-return-field} to themselves. Consequently $\iota$, $\tau_1$, $\mathrm{pr}$, $\psi_1^{-1}$ and
$\mathcal T_{\rm u}$ commute with $\mathsf r$.
\end{lemma}

\begin{proof}
$\mathsf r$ is unitary on $\C^3$ and $\mathcal F\circ\mathsf r=\mathcal F$, so $\mathsf r$ preserves $\|Z\|$, $\omega_{\rm FS}$, $\|W\|$, $\Omega_E$, $X_T$, $u$, $p$,
$A$, $B$, $A_s$, $B_s$. Since $\bar\zeta^2=\zeta$, $n(\mathsf r Z)=3\,\overline{\mathsf r Z}^2-3\bar f\,\mathsf r Z=\mathsf r\,n(Z)$; so $G$ and $N_{\rm den}$ are invariant,
and both fields satisfy $\dot Z(\mathsf r Z,W)=\mathsf r\,\dot Z(Z,W)$, $\dot W(\mathsf r Z,W)=\dot W(Z,W)$. By uniqueness of solutions the flows commute
with $\mathsf r$. $\iota$ commutes with $\mathsf r$ because $p$ is invariant, and $\mathrm{pr}$ commutes with $\mathsf r$ because the lift of $\mathsf r$ to
$E_3$ covers $\mathsf r$.
\end{proof}

\begin{lemma}[Injectivity modulo gauge]\label{lem:rtc-injective}
\lean{SymplecticPacking.Return.compositeT_inj_mod_gauge}
For $x,x'\in\mathcal N_h$ and $|c|=1$, the equality
$\mathcal T_{\rm u}(x)=c\mathcal T_{\rm u}(x')$ implies
$x=(cZ',c^3U')$.
\end{lemma}
\begin{proof}
Each factor of $\mathcal T_{\rm u}=\psi_1^{-1}\circ\mathrm{pr}\circ\tau_1\circ\iota$ is injective on representatives: $\iota$ has the
explicit inverse of Lemma~\ref{lem:rtc-injection}, $\tau_1$ and $\psi_1^{-1}$ are time maps of complete flows, and $\mathrm{pr}$ is
injective on $X_1$ because $W=\mathcal F(Z)$ there. Each factor commutes with $(Z,W)\mapsto(cZ,c^3W)$, $|c|=1$, so
$\mathcal T_{\rm u}(cZ',c^3U')=c\,\mathcal T_{\rm u}(Z',U')$. The hypothesis therefore reads
$\mathcal T_{\rm u}(x)=\mathcal T_{\rm u}(cZ',c^3U')$, and injectivity gives $x=(cZ',c^3U')$.
\end{proof}


\begin{proposition}[Avoidance and compact margin]\label{prop:rtc-avoid}
\lean{SymplecticPacking.Return.avoid_L,SymplecticPacking.Return.avoid_L_compact}
If $x=(Z,U)\in\mathcal N_h$ and $Z_0\ne0$, then
$(\mathcal T_{\rm u}(x))_0\ne0$ and its last two coordinates do not vanish together.
On every nonempty compact family of these inputs,
$m_{\mathcal K}=\min_{x\in\mathcal K}|(\mathcal T_{\rm u}(x))_0|^2>0$. For $\mathcal K=\varnothing$ set $m_{\mathcal K}=1$.
\end{proposition}
\begin{proof}
Let $Y=\mathcal T_{\rm u}(x)$, a unit vector. If $Y_0=0$, then $\mathsf r Y=Y$, so
$\mathcal T_{\rm u}(\mathsf r x)=\mathsf r\mathcal T_{\rm u}(x)=Y=\mathcal T_{\rm u}(x)$ by Lemma~\ref{lem:rtc-equivariant}, and
Lemma~\ref{lem:rtc-injective} with $c=1$ gives $\mathsf r x=x$; by Lemma~\ref{lem:rtc-fixed}, $Z_0=0$. If $Y_1=Y_2=0$, then $\mathsf r Y=\zeta Y$, so
$\mathcal T_{\rm u}(\mathsf r x)=\zeta\mathcal T_{\rm u}(x)$ and Lemma~\ref{lem:rtc-injective} gives $\mathsf r x=(\zeta Z,\zeta^3U)$; by
Lemma~\ref{lem:rtc-fixed}, again $Z_0=0$. Both contradict $Z_0\ne0$. The function $x\mapsto|(\mathcal T_{\rm u}(x))_0|^2$ is continuous
and positive on $\mathcal K$, so it has a positive minimum when $\mathcal K\ne\varnothing$.
\end{proof}


\begin{lemma}[The ball chart on representatives]\label{lem:rtc-ball-chart}
\lean[needed]{SymplecticPacking.Return.ballChart, SymplecticPacking.Return.ballChart_gauge,
 SymplecticPacking.Return.ballChart_symplectic_rep, SymplecticPacking.Return.ballChart_affine,
 SymplecticPacking.Return.ballChart_norm, SymplecticPacking.Return.ballAffine_eq, SymplecticPacking.Return.ballChart_norm_lt}
On $S^5_*=\{Z\in\C^3:|Z|=1,\ Z_0\ne0\}$ put $b_{\rm ball}(Z)=\frac{\bar Z_0}{|Z_0|}\frac{(Z_1,Z_2)}{\sqrt\pi}$. This map is gauge
invariant, $\pi|b_{\rm ball}(Z)|^2=1-|Z_0|^2<1$, and in the affine coordinates $z=(Z_1/Z_0,Z_2/Z_0)$ it is the chart
$\mathrm{ch}(z)=z/\sqrt{\pi(1+|z|^2)}$ of Lemma~\ref{lem:pre-fs}(b). Moreover
$b_{\rm ball}^*\omega_0=\frac1\pi\sum_{j=0}^2dx_j\wedge dy_j\big|_{S^5_*}=\omega_{\rm FS}|_{S^5_*}$, where $Z_j=x_j+iy_j$. Hence
$(b_{\rm ball}\circ\varphi)^*\omega_0=\varphi^*\omega_{\rm FS}$ for every smooth map $\varphi$ of a manifold into $S^5_*$.
\end{lemma}

\begin{proof}
The unit phases cancel, and substituting $Z=e^{i\chi}(1,z)/\sqrt{1+|z|^2}$ gives the affine formula and the norm identity.
Write $e=\bar Z_0/|Z_0|=e^{i\phi}$, $\phi=-\arg Z_0$, and $y_j=eZ_j/\sqrt\pi$ ($j=1,2$). Then
$\bar y_jdy_j=\frac1\pi\big(|Z_j|^2\,\bar e\,de+\bar Z_jdZ_j\big)$ with $\bar e\,de=i\,d\phi$, so the Liouville form
$\lambda_0=\frac12\sum_{j=1}^2\operatorname{Im}(\bar y_jdy_j)$ pulls back to
$\frac1{2\pi}\big((1-|Z_0|^2)d\phi+\sum_{j=1}^2\operatorname{Im}(\bar Z_jdZ_j)\big)$ on $S^5_*$. Since
$\operatorname{Im}(\bar Z_0dZ_0)=|Z_0|^2d\arg Z_0=-|Z_0|^2d\phi$, this equals $\frac1{2\pi}\big(d\phi+\operatorname{Im}(Z^\dagger dZ)\big)$, whose
differential is $\frac1{2\pi}\sum_{j=0}^2d(x_jdy_j-y_jdx_j)=\frac1\pi\sum_jdx_j\wedge dy_j$. As $d\lambda_0=\omega_0$, this proves the first
identity; the second holds because $\frac1{2\pi}\operatorname{Im}(Z^\dagger dZ)=\eta_H$ on the sphere and $d\eta_H=\omega_{\rm FS}$.
\end{proof}

The vector
$(\sqrt{1-\pi|y|^2},\sqrt\pi\,y_1,\sqrt\pi\,y_2)$ is a unit vector with positive first coordinate for $y\in\Bo(1)$,
and its ball-chart image is $y$. Conversely, every unit vector with nonzero first coordinate has a unique
unit-gauge representative with positive first coordinate, namely the vector above for its ball-chart image.

\subsection{The transfer theorem}\label{sec:rtc-transfer}

\begin{theorem}[Transfer on the cover]\label{thm:rtc-transfer}\label{thm:cub-transfer}
\lean[needed]{SymplecticPacking.Return.compositeT_inj_mod_gauge,SymplecticPacking.Return.compositeT_pullback,SymplecticPacking.Return.compositeT_equivariant,SymplecticPacking.Return.avoid_L_compact,
SymplecticPacking.Return.injection_pullback, SymplecticPacking.Return.ballChart_symplectic_rep,
SymplecticPacking.Embed.Cubic.fcap_eq_imp, SymplecticPacking.Embed.Cubic.fcap_eq_imp}
\uses{cons:cub-transfer, prop:rtc-central, prop:rtc-transport, prop:rtc-return, prop:rtc-avoid, lem:rtc-ball-chart}
Let $w_*=h/(d_n-h)$ and $\mathcal T_{\rm u}=\psi_1^{-1}\circ\mathrm{pr}\circ\tau_1\circ\iota$ on $\mathcal N_h$. The cover map
\[
 \mathcal T=b_{\rm ball}\circ\mathcal T_{\rm u}\circ\mathsf N
 \quad\text{on}\quad\{p<h,\ (u,v)\notin\mathcal Z\}
\]
is smooth, invariant under the deck maps, takes values in $\Bo(1)$ and satisfies $\mathcal T^*\omega_0=\Omega_0$. Two points
have the same image if and only if they are deck equivalent. For every compact $\mathcal K$ in the domain,
$\mathcal T(\mathcal K)\subset B(1-m_{\mathcal K})$ with $m_{\mathcal K}>0$ as in Proposition~\ref{prop:rtc-avoid}. In particular the
restriction of $\mathcal T$ to an open set on which it is injective is a symplectic embedding, and $\mathcal T$ carries
pairwise disjoint subsets of its domain, no two points of whose union are deck equivalent, to pairwise disjoint sets.
\end{theorem}
\begin{proof}
On $\{p<h\}$, $\mathsf N$ takes values in $\mathcal N_h$ with $p<d_n$, and $\tau_1\circ\iota$ lands in
$X_1\cap\{|W|^2\le w_*\}$ (Lemma~\ref{lem:rtc-injection}, Proposition~\ref{prop:rtc-transport}). Composing
Propositions~\ref{prop:rtc-central}, \ref{prop:rtc-transport} and~\ref{prop:rtc-return} gives
$(\mathcal T_{\rm u}\circ\mathsf N)^*\omega_{\rm FS}=\Omega_0$. If $(u,v)\notin\mathcal Z$, then $Z_0(u,v)\ne0$, so
$\mathcal T_{\rm u}(\mathsf N(x))\in S^5_*$ by Proposition~\ref{prop:rtc-avoid}; Lemma~\ref{lem:rtc-ball-chart} then gives
$\mathcal T^*\omega_0=\Omega_0$ and $\pi|\mathcal T|^2=1-|(\mathcal T_{\rm u}\mathsf N)_0|^2<1$. A deck map changes $\mathsf N(x)$ to a
gauge-equivalent pair (Construction~\ref{cons:fer-atlas}(b)), $\mathcal T_{\rm u}$ is gauge equivariant and $b_{\rm ball}$ is gauge
invariant, so $\mathcal T$ is deck invariant. Conversely, suppose $\mathcal T(x)=\mathcal T(x')$. A unit vector with nonzero first
coordinate is determined up to gauge by its ball-chart image (see the paragraph before this subsection), so
$\mathcal T_{\rm u}(\mathsf N x)=c\,\mathcal T_{\rm u}(\mathsf Nx')$ with $|c|=1$, and Lemma~\ref{lem:rtc-injective} gives
$\mathsf N(x)\sim\mathsf N(x')$. Then $[Z(u,v)]=[Z(u',v')]$ in $\CP^2$, and injectivity of $\Phi_\Gamma$ (Lemma~\ref{lem:fer-torus})
gives $(u,v)-(u',v')\in\Z^2$. Applying the corresponding deck map $T$ to $x'$ we get $\mathsf N(x)\sim\mathsf N(Tx')$ with equal unit
vectors $Z(u,v)$, so the gauge factor is one and the fibre coordinates agree: $x=Tx'$. The margin is
$1-\pi|\mathcal T|^2=|(\mathcal T_{\rm u}\mathsf N)_0|^2\ge m_{\mathcal K}$. Finally, on an open set on which $\mathcal T$ is injective it
is an injective immersion between manifolds of equal dimension, hence an open symplectic embedding, and the last
assertion follows from the characterisation of equal images.
\end{proof}

All objects in the theorem are given by closed formulas and by the solutions of the explicit ODEs \eqref{eq:rtc-transport}
and \eqref{eq:rtc-return-field}, the latter integrated backwards from $s=1$ to give $\psi_1^{-1}$.

%----------------------------------------------------------------------
\subsection{The packing in dimension four}\label{sec:rtc-packing}

We now compose the packing in the model with the transfer. Recall from Lemma~\ref{lem:cub-auxiliary} the data attached to
capacities $0<\rho'_1,\dots,\rho'_k<\frac13$ with $\sum_i(\rho'_i)^2<1$: auxiliary capacities $a_i$, a height $h$, the core fraction
$\tau_\Delta=(a_{\max}+h)/(2h)$, and $\Pb(p)=\Hb(p)-\sum_i(a_i-p)_+$ on $\Delta=[0,h]$. The transfer depends on the capacities only
through $h$.

\begin{construction}[Data of the four-dimensional packing]\label{cons:rtc-data}
\uses{lem:cub-auxiliary, M-prop:cmp-linear, prop:scalar-cubic-data, def:shd4-data, M-thm:mech-packing, prop:cub-handle,
 prop:fer-handle, thm:rtc-transfer, prop:rtc-avoid}
Let $k\ge1$ and $0<\rho_1,\dots,\rho_k<\frac13$ with $\sum_i\rho_i^2<1$. Define $(a_i,h,P,(K,q),\gamma)$ as follows.
\begin{enumerate}[label=(\alph*)]
\item\emph{Equal case}: $k\ge10$, $\rho_1=\dots=\rho_k=\rho$ and $\frac34<\rho\sqrt k<1$. Put $t_0=\frac12(1+\rho\sqrt k)$, $a_i=a=t_0/\sqrt k$, $h=t_0/3$,
$P=\Pb$, and let $(K,q)$ be the comparison datum of \Mref{Proposition}{prop:cmp-linear}, with gap
$\gamma=\frac{31}{32}\cdot3(1-t_0)$.
\item\emph{General case}, otherwise. Use the auxiliary quantities of
\eqref{eq:scalar-cubic-parameters} for $\rho$. Define $K$ by the scalar iterate and fixed convolution
\eqref{eq:scalar-iteration}--\eqref{eq:scalar-mollifier}, and use the planar path
\eqref{eq:scalar-alexander}--\eqref{eq:scalar-qhat}. The cubic comparison data are
\eqref{eq:scalar-cubic-planar}, with $g=3d/4$ and comparison gap $g$ as in
Proposition~\ref{prop:scalar-cubic-data}.
\end{enumerate}
Every admissible input falls under (a) or (b); in particular equal capacities with $k\le9$, or with $\rho\sqrt k\le\frac34$, are
treated by (b). (For $k=9$, the formulas of (a) would give $a=t_0/3=h$, violating $a<h$.)

In both cases put $d_n=\frac12(h+\frac13)$, let $f_i=f_i^{\rm cap}$ be the maps of Remark~\ref{rem:shd-four}, built by the handle construction of
\S\ref{sec:shd-four}, put $\bar\rho_i=(\rho_i+a_i)/2$ and $\hat\rho_i=\bar\rho_i+(a_i-\bar\rho_i)/4$, and let the open domain be
$\Bo(\hat\rho_i)$. Let $\rho_i^+<\rho_i^{\rm dom}$ be the recorded capacities of Remark~\ref{rem:rtc-recorded}. Put
$\mathcal K=\bigcup_i f_i(B(\rho_i^+))$ and $m=\frac12\min_{z\in\mathcal K}(1-\pi|z|^2)$.
\end{construction}

\begin{theorem}[Packing in dimension four]\label{thm:cub-packing}\label{prop:rtc-embeddings}
\lean[needed]{SymplecticPacking.Branch2.embeddings2_spec,
SymplecticPacking.Branch2.CubicData.Valid.fcapGen_injective, SymplecticPacking.Embed.Cubic.fcapV_disjoint,
SymplecticPacking.Embed.Cubic.fcapV_mem_ball, SymplecticPacking.Embed.Cubic.fcap_injective, SymplecticPacking.Embed.Cubic.fcap_disjoint,
SymplecticPacking.Embed.Cubic.fcap_mem_ball, SymplecticPacking.Embed.Cubic.handle_phiInvSpec,
SymplecticPacking.Handle.Cubic.handle_deckEquivariant}
\uses{cons:rtc-data, lem:cub-auxiliary, prop:scalar-cubic-data, prop:scalar-planar-realisation, M-prop:cmp-linear,
 def:shd4-data, M-thm:mech-packing, prop:cub-handle, prop:fer-handle, M-prop:mech-assembly, thm:rtc-transfer, prop:rtc-avoid}
Let $k\ge1$ and $0<\rho_1,\dots,\rho_k<\frac13$ with $\sum_i\rho_i^2<1$. Then there are explicit numbers
$\rho_i<\rho_i^+<\rho_i^{\rm dom}$ and smooth maps $f_i\colon\Bo(\rho_i^{\rm dom})\to\C^2$ with symplectic derivative, injective
on $B(\rho_i^+)$, such that the sets $f_i(B(\rho_i^+))$ are pairwise disjoint and contained in $B(1-\mu_{\rm ball})$ for some
$\mu_{\rm ball}\in(0,1)$. More precisely, for the data of Construction~\ref{cons:rtc-data}:
\begin{enumerate}[label=(\roman*)]
\item each $f_i$ is a smooth symplectic embedding of $\Bo(\hat\rho_i)$ into $(\Bo(1),\omega_0)$, and
$\rho_i<\rho_i^+<\rho_i^{\rm dom}\le\hat\rho_i$;
\item the sets $f_i(\Bo(\hat\rho_i))$, $1\le i\le k$, are pairwise disjoint;
\item $m>0$ and $f_i(B(\rho_i^+))\subset B(1-m)$ for every $i$;
\item the comparison isotopy $q$, the Moser map $\Phi$, the handle map $\mathcal D$ and the transfer $\mathcal T$ are the same for all $i$;
only the layer map $E_i$ depends on $i$.
\end{enumerate}
\end{theorem}

\begin{proof}
In both cases the data $(k,\rho_i,a_i,h,d_n,g,K,q)$ are valid cubic data (Definition~\ref{def:shd4-data}): in case (a) by
\Mref{Proposition}{prop:cmp-linear} and Remark~\ref{rem:rtc-equal}, in case (b) by Proposition~\ref{prop:scalar-cubic-data};
moreover $h<d_n=\frac12(h+\frac13)<\frac13$. The maps are $f_i=\mathcal T\circ f^N_i$ with $f^N_i$ of Proposition~\ref{prop:shd4-maps}
and $\mathcal T$ of Theorem~\ref{thm:rtc-transfer} (Remark~\ref{rem:shd-four}).

(i) By Proposition~\ref{prop:shd4-maps}(a), $f^N_i$ is smooth on $\Bo(\hat\rho_i)$ with values in $\{p<h\}\setminus\widetilde{\mathcal F}$, the
domain of $\mathcal T$, and $(f^N_i)^*\Omega_0=\omega_0$; by Theorem~\ref{thm:rtc-transfer}, $f_i$ is smooth with values in $\Bo(1)$
and $f_i^*\omega_0=\omega_0$. If $f_i(X)=f_i(X')$, Theorem~\ref{thm:rtc-transfer} gives $f^N_i(X)=Tf^N_i(X')$ for a deck map $T$, and
Proposition~\ref{prop:shd4-maps}(c) gives $X=X'$. A symplectic map is an immersion, and an injective immersion between
manifolds of equal dimension is an open embedding; so $f_i$ is a symplectic embedding of $\Bo(\hat\rho_i)$. The inequalities
for the recorded capacities are in Remark~\ref{rem:rtc-recorded}.

(ii) If $f_i(X)=f_j(X')$, the same argument with Proposition~\ref{prop:shd4-maps}(c) gives $i=j$ and $X=X'$.

(iii) $\mathcal K$ is a finite union of continuous images of compact balls, hence compact, and lies in $\Bo(1)$; so the continuous
positive function $1-\pi|z|^2$ has a positive minimum on it, $m>0$, and $1-\pi|z|^2\ge2m>m$ on $\mathcal K$.

(iv) The factors $q$ (in $\mathcal D$), $\mathcal D$, $\Phi$, $\chi_{\mathcal A}$ and $\mathcal T$ do not depend on $i$.

The first statement follows from (i)--(iii) with $\mu_{\rm ball}=m<1$, restricting $f_i$ to $\Bo(\rho_i^{\rm dom})$.
\end{proof}

\begin{remark}[Recorded capacities]\label{rem:rtc-recorded}
The proof gives embeddings of the open balls $\Bo(\hat\rho_i)$ with pairwise disjoint images, so Theorem~\ref{thm:cub-packing}
holds for every choice of $\rho_i<\rho_i^+<\rho_i^{\rm dom}\le\hat\rho_i$. The choice enters only the radii of the common
isotopy (\Mref{Definition}{def:ham-radii}) and hence the Hamiltonian. The construction records
$(\rho_i^+,\rho_i^{\rm dom})=(\bar\rho_i,\hat\rho_i)$ in case (b), as in \eqref{eq:scalar-enlargement}, and
$(\rho_i^+,\rho_i^{\rm dom})=(\rho r_3^2,\rho r_4^2)$ in case (a), where $r_4=\sqrt{\bar\rho/\rho}$ and $r_3=(1+r_4)/2$. In case (a),
$1<r_3<r_4$ gives $\rho<\rho r_3^2<\rho r_4^2=\bar\rho<\hat\rho$. For target capacity $C$ and given capacity $c=C\rho$
the physical map $u\mapsto\sqrt C\,f_i(u/\sqrt C)$ has the capacities $C\rho_i^+$ and $C\rho_i^{\rm dom}$, that is, $cr_3^2$
and $cr_4^2$ in case (a).
\end{remark}

\begin{remark}[The equal-capacity branch]\label{rem:rtc-equal}\label{rem:cub-lean}
\lean{SymplecticPacking.Params.t0, SymplecticPacking.Params.a, SymplecticPacking.Params.h,
 SymplecticPacking.Params.dn_sub_h, SymplecticPacking.Params.one_sub_three_dn, SymplecticPacking.Params.a_sub_rho,
 SymplecticPacking.Params.density, SymplecticPacking.Branch2.CubicData.ofParams_valid,
 SymplecticPacking.Branch2.embeddingsEq, SymplecticPacking.Branch2.embeddingsEq_spec,
 SymplecticPacking.Branch2.cubicEmbeddings_spec}
\uses{cons:rtc-data, thm:cub-packing, rem:fer-lean}
In case (a) of Construction~\ref{cons:rtc-data} let $\eta_\star=1-\rho\sqrt k\in(0,\frac14)$ be the volume defect. Then
\[
 t_0=1-\frac{\eta_\star}2,\quad a=\frac{t_0}{\sqrt k},\quad h=\frac{t_0}3,\quad d_n=\frac13-\frac{\eta_\star}{12},\qquad
 a-\rho=\frac{\eta_\star}{2\sqrt k},\quad d_n-h=\frac{\eta_\star}{12},\quad 1-3d_n=\frac{\eta_\star}4,
\]
the comparison function is $K(p)=(3\sqrt k-9)p$, and the comparison gap of \Mref{Proposition}{prop:cmp-linear} is
$\frac{31}{32}g$ with $g=3(1-t_0)=\frac{3\eta_\star}2$. With the isotopy of Theorem~\ref{thm:pll-inner} and the handle
construction of \S\ref{sec:shd-four}, the maps of this branch are
\[
 f_i= b_{\rm ball}\circ\psi_1^{-1}\circ\operatorname{pr}\circ\tau_1\circ\iota
       \circ\mathsf N\circ\Phi^{-1}\circ\chi_{\mathcal A}\circ\mathcal D^{-1}\circ E_i\qquad(1\le i\le k).
\]
For the input of \Mref{Corollary}{thm:main-epoch}, $\rho=\sqrt{(1-\eps/2)/k}$ and
$\eta_\star=1-\sqrt{1-\eps/2}\in(\eps/4,\eps/2)$; case (a) is selected exactly when $k\ge10$ and $\rho\sqrt k>3/4$,
and case (b) covers the other inputs. For unit balls (capacity $\pi$) the target capacity is $C=\pi/\rho$, and the
physical rescaling of \Mref{Definition}{def:ham-normalise} makes the target the open Euclidean ball of radius
$R=\rho^{-1/2}$, so $k/R^4=k\rho^2=(1-\eta_\star)^2$. For the data of \Mref{Theorem}{thm:intro-4d}, $R^4=k/(1-\eps/4)$, so
$\rho\sqrt k=\sqrt{1-\eps/4}>\frac34$ and $\eta_\star=1-\sqrt{1-\eps/4}\in(\eps/8,\eps/4)$; case (a) applies for every
$\eps\in(0,1)$ when $k\ge10$. The strip constants are those of Remark~\ref{rem:shd-four} with
$\delta=(a-\rho)/2$.

The defect $1-\rho\sqrt k=\eta_\star$ splits as $(1-t_0)+\sqrt k\,(a-\rho)=\frac{\eta_\star}2+\frac{\eta_\star}2$. The fibre part $1-t_0=1-3h$ pays for
$h<\frac13$, for the comparison gap $g$ and for the margins $d_n-h$ and $1-3d_n$ of the transfer; the strip part $\sqrt k\,(a-\rho)$ pays for the enlargement from $\rho$ to $\bar\rho$ and for the excess of the strip map.
The transfer, the ball chart and the dilation are symplectic embeddings and lose no volume beyond these margins.
\end{remark}

\begin{remark}[Obstructions]\label{rem:cub-obstructions}
The condition $3\rho_i<1$ of Theorem~\ref{thm:cub-packing} is not necessary (one ball of capacity $\rho<1$ embeds by
inclusion), but it cannot be replaced by the volume and two-ball inequalities: five balls of capacity
$\frac{11}{25}>\frac13$ and five of capacity $\frac1{100}$ satisfy $\sum_i\rho_i^2=\frac{1937}{2000}<1$ and
$\rho_i+\rho_j\le\frac{22}{25}<1$, yet by the disjoint-union property of ECH capacities \cite{Hutchings2011} the fifth ECH
capacity of their union is at least $5\cdot\frac{11}{25}=\frac{11}5$, while that of $\Bo(1)$ is $2$; so they do not embed.
\end{remark}
